% Agent 03: Pages 352--355 (book pp.352--355, PDF pp.5--7)
% Content: End of Section 2.1 (from eq.~2.1.20 onward), Figure 2.1.2,
%          beginning of Section 2.2 (Free-Bound Radiation)
% Equations: (2.1.20)--(2.1.31), (2.2.1)--(2.2.9)

%%---------------------------------------------------------------------------
%% Continuation of Section 2.1
%% Small-angle, uncertainty-principle region (cont.)
%%---------------------------------------------------------------------------

Thus, the uncertainty principle (``$\geq$'' in above equation) prevents the existence
of scatterings with $b$ less than the electron de~Broglie wavelength
\begin{equation}
b_{dB} = \hbar / m_e v.
\tag{2.1.20}
\end{equation}

If $b_{dB} < b_{s-l}$, this limitation has no effect on small-angle scatterings; and the
small-angle, classical Gaunt factor (2.1.19) is valid.  If $b_{dB} > b_{s-l}$, the uncertainty
principle comes into play, and one must use $b_{dB}$ rather than $b_{s-l}$ as the lower
limit on the small-angle integral (2.1.17) for $d\sigma/dv$.  The result is
\begin{equation}
\frac{d\sigma}{dv} = \frac{16\,\alpha^2}{3}\,\frac{Z^2 r_0^2}{v^2}\,
\ln\!\left(\frac{b_{\max}}{b_{dB}}\right).
\tag{2.1.21}
\end{equation}

A more exact, quantum mechanical calculation gives a slightly different argument
in the logarithm,
\begin{equation}
\frac{d\sigma}{dv} = \frac{16\,\alpha^2}{3}\,\frac{Z^2 r_0^2}{v^2}\,
\ln\!\left(\frac{2\,b_{\max}}{b_{dB}}\right).
\tag{2.1.22}
\end{equation}

corresponding to the Gaunt factor
\begin{equation}
G(\nu,v) = \frac{\sqrt{3}}{\pi}\,
\ln\!\left(\frac{4\,\frac{1}{2}m_e v^2}{h\nu}\right).
\tag{2.1.23}
\end{equation}

Thus, in the large-angle region $G = 1$; and in the forbidden region $G = 0$.  It turns
out that throughout the photon-discreteness portion of the large-angle region,
$G$ remains approximately 1; see Figure~2.1.1.

\textit{Small-angle, classical region}\quad In the region $\tau = 1/\omega \gg b_{s-l}/v$---i.e.
\begin{equation}
\frac{h\nu}{\frac{1}{2}m_e v^2}
\ll \left(\frac{\frac{1}{2}m_e v^2}{Z^2 R_y}\right)^{1/2}
\tag{2.1.14}
\end{equation}

In the ``photon-discreteness'' portion of the small-angle, uncertainty-principle
region, discreteness effects modify the Gaunt factor (2.1.22) into the form
\begin{subequations}
\begin{align}
G(\nu,v) &= \frac{\sqrt{3}}{\pi}\,
\ln\!\left[\frac{\frac{1}{2}m_e(v_i + v_f)^2}{h\nu}\right],
\tag{2.1.24a}\\[6pt]
G(\nu,v) &= \frac{\sqrt{3}}{\pi}\,
\ln\!\left[\frac{v_i + v_f}{v_i - v_f}\right]^2,
\tag{2.1.24b}
\end{align}
\end{subequations}

Here $v_i$ and $v_f$ are the electron speeds before and after the collision:
\begin{equation}
\tfrac{1}{2}m_e v_f^2 = \tfrac{1}{2}m_e v_i^2 - h\nu.
\end{equation}

[Expression (2.1.24) is the result of a quantum-mechanical Born-approximation
calculation, valid throughout the small-angle, uncertainty-principle region.]

%%---------------------------------------------------------------------------
%% Radiation from a plasma
%%---------------------------------------------------------------------------

\bigskip
\textit{Radiation from a plasma}

Consider a plasma which may contain a variety of ionic species.  Focus attention
on radiation from all ions of a particular type, with charge $Ze$ and atomic weight
$A$.  Let $C_A$ be the concentration by mass of the species of atom that gives rise to
this ion, and let $f_i$ be the fraction of all such atoms ionized to the state of interest.
Let $f_e$ be the number of unbound electrons per baryon in the gas.  Then
\begin{equation}
n_I = \text{(number of ions per unit volume)} = f_i C_A \rho_0 / A m_p.
\tag{2.1.25}
\end{equation}

Here $\rho_0$ is the density of rest mass.  (Of course, $f_i \leq 1$; $C_A \leq 1$.)  The ions
and electrons both have Maxwell velocity distributions; but because $m_e \ll A m_p$,
the velocities of the electrons and their accelerations during Coulomb scattering
are far greater than those of the ions.  Therefore, in calculating bremsstrahlung
one can regard the ions as at rest.  The number of photons per unit mass of plasma
is $f_e/m_p$, so the total power per unit frequency emitted from one gram of plasma
by electron-ion scatterings is
\begin{equation}
\varepsilon_\nu
= \frac{C_A f_i}{A m_p}
\underbrace{\vphantom{\bigg|}\quad}_{\substack{
\text{fractions of all} \\
\text{ions per} \\
\text{unit mass}
}}
\cdot \int
\underbrace{\vphantom{\bigg|}\quad}_{\substack{
\text{number} \\
\text{of electrons}
}}
\underbrace{\vphantom{\bigg|}\quad}_{\substack{
\text{electrons that have} \\
\text{speeds $v$ in range $dv$}
}}
\frac{f_e \rho_0}{m_p}\;v\,h\nu\;
\frac{d\sigma}{dv}\;dv.
\tag{2.1.26}
\end{equation}

Here $T$ is the temperature in the plasma.  By inserting into the integrand
expressions (2.1.12) and (2.1.13) for $d\sigma/dv$ and performing the integration, one
obtains
\begin{align}
\varepsilon_\nu
&= \frac{C_A f_i f_e^2}{A m_p}\,
\frac{32\pi e^6}{3\,m_e c^3}\,
\left(\frac{2\pi m_e}{3kT}\right)^{\!1/2}\!
\frac{\rho_0^2}{m_p^2}\;
e^{-h\nu/kT}\;
\bar{G}(\nu,T)
\notag\\[4pt]
&= 2.5 \times 10^{10}\;
\frac{\text{erg}}{\text{g\;sec\;Hz}}\;
\left(\frac{C_A f_i f_e^2}{A}\right)
\left(\frac{\rho_0}{\text{g\;cm}^{-3}}\right)
\left(\frac{\rho_0}{A}\right)
T_K^{\,1/2}\;e^{-h\nu/kT}\;\bar{G}(\nu,T).
\tag{2.1.27}
\end{align}

%%---------------------------------------------------------------------------
%% Mean Gaunt factor \bar{G}
%%---------------------------------------------------------------------------

Here $T_K$ is temperature measured in ${}^\circ$K, and $\bar{G}(\nu,T)$ is the ``Maxwell--Boltzmann-averaged
Gaunt factor'', also called the ``mean Gaunt factor'':
\begin{equation}
\bar{G}(\nu,T) = \int_0^\infty
G\!\left(\nu,\;v = \left[\frac{2(xkT + h\nu)}{m_e}\right]^{1/2}\right)
e^{-x}\,dx.
\tag{2.1.28}
\end{equation}

Corresponding to the 2-dimensional plot of Gaunt factors (Figure~2.1.1) one
can make a 2-dimensional plot of mean Gaunt factors (Figure~2.1.2).  It is
straightforward to calculate the mean Gaunt factors shown there (up to
quantities of order unity) from the Gaunt factors of Figure~2.1.1.  Notice that,

%%---------------------------------------------------------------------------
%% Figure 2.1.2 — Mean Gaunt factors
%%---------------------------------------------------------------------------

\begin{figure}[htbp]
\centering
% \includegraphics[width=\columnwidth]{fig_2_1_2}
\fbox{\parbox{0.9\columnwidth}{\centering [Figure 2.1.2 placeholder]\\[6pt]
Mean Gaunt factors for bremsstrahlung of frequency $\nu$ emitted by
electron-ion collisions in a plasma of temperature $T$.\\[4pt]
Horizontal axis: $h\nu/kT$.\\
Vertical axis: $(kT)/(Z^2 R_y)$.\\[4pt]
Regions shown:\\
--- ``Small-angle, u.p., tail region'',
$\bar{G} = \bigl(\frac{3}{\pi}\bigr)\,\bigl(\frac{kT}{h\nu}\bigr)^{1/2}$\\
--- ``Small-angle, u.p.\ region'',
$\bar{G} = \frac{\sqrt{3}}{\pi}\ln\!\left[\frac{4}{\xi}\,\frac{kT}{h\nu}\right]$\\
--- ``Small-angle, classical region'',
$\bar{G} = \frac{\sqrt{3}}{\pi}\,\ln\!\left[\left(\frac{4}{\xi^{3/2}}\right)
\left(\frac{kT}{h\nu}\right)\left(\frac{kT}{Z^2 R_y}\right)^{1/2}\right]$\\
--- ``Large-angle region'', $\bar{G} = 1$\\
--- ``Large-angle, tail region'', $\bar{G} = 1$\\
--- ``Forbidden region''}}
\caption{\textit{Figure~2.1.2.}\ Mean Gaunt factors for bremsstrahlung of frequency $\nu$ emitted by
electron-ion collisions in a plasma of temperature~$T$.}
\label{fig:2.1.2}
\end{figure}

aside from quantities of order unity, the mean Gaunt factors in the large-angle
and small-angle regions are obtained by simply taking the Gaunt
factors and making the replacement $\frac{1}{2}m_e v^2 \approx kT$:
\begin{equation}
\bar{G}(\nu,kT) \simeq G\bigl(\nu,\;\tfrac{1}{2}m_e v^2 = kT\bigr).
\end{equation}

Notice also that the ``forbidden region'' of the Gaunt-factor plot is replaced, in
the mean-Gaunt-factor plot, by ``tail regions''.  The radiation in these tail regions
is produced by the high-energy tail, $\frac{1}{2}m_e v^2 \approx h\nu \gg kT$, of the Maxwell--Boltzmann
distribution.

Tables of mean Gaunt factors, calculated with much higher accuracy than
Figure~2.1.2, are given by \citet{Green1959} and by \citet{Karzas1961}.

From Figure~2.1.1 it is clear that the uncertainty principle cannot have any
effect on the large-angle region.

%%---------------------------------------------------------------------------
%% Total emissivity and frequency-averaged mean Gaunt factor
%%---------------------------------------------------------------------------

The total energy emitted by electron-ion collisions is obtained by integrating
expression (2.1.27) over frequency:
\begin{align}
e &= \int \varepsilon_\nu\,d\nu
= \frac{16\,C_A f_i f_e^2}{3}\;
\frac{\alpha\,m_e c^2_0}{A}\;
\left(\frac{2\pi kT}{m_e c^2}\right)^{\!1/2}\!
\frac{\rho_0}{m_p^2}\;
\left(\frac{C_A f_i f_e^2}{A}\right)
\left(\frac{\rho_0}{\text{g\;cm}^{-3}}\right)
T_K^{\,1/2}\;\bar{G}(T)
\notag\\[4pt]
&= 5.2 \times 10^{20}\;
\frac{\text{ergs}}{\text{g\;sec}}\;
\left(\frac{C_A f_i f_e^2}{A}\right)
\left(\frac{\rho_0}{\text{g\;cm}^{-3}}\right)
T_K^{\,1/2}\;\bar{G}(T).
\tag{2.1.29}
\end{align}

Here $\bar{G}(T)$ is the frequency-averaged mean Gaunt factor
\begin{equation}
\bar{G}(T) = \int_0^\infty \bar{G}(\nu = xkT/h,\;T)\;e^{-x}\,dx.
\tag{2.1.30}
\end{equation}

From a value of $\bar{G} \simeq 1.2$ at $kT/Z^2 R_y \sim 0.01$, it increases gradually to a maximum
of $1.42$ at $kT/Z^2 R_y \sim 1$, and then decreases gradually toward an asymptotic,
high-temperature value of $1.103$ at $kT/Z^2 R_y > 100$; see Figure~2 of \citet{Green1959}.

In a plasma at nonrelativistic temperatures, the radiation from electron-electron
collisions and from ion-ion collisions is negligible compared to electron-ion
radiation.  This is because (i) when identical particles scatter, the first time
derivative of the electric dipole moment is conserved, so only radiation of
quadrupole order and higher [with intensity $\sim(v/c)^2$ less than dipole radiation]
can be emitted; (ii) the speeds and accelerations of colliding ions are far less than
those of electrons because of their much larger masses.

Equations (2.1.27) and (2.1.29), together with Figures~2.1.2 and~2.1.3, are
the chief results of astrophysical interest from this section.  When dealing with
atoms that are only partially ionized, one must sometimes correct these results
for screening of the nuclear charge by the bound electrons; see, e.g.,
\citet{Ginzburg1967}.  For temperatures approaching and exceeding $kT = m_e c^2$,
relativistic effects and electron-electron collisions modify the emissivity
(2.1.29) to the form
\begin{align}
e &= 5.2 \times 10^{20}\;
\frac{\text{ergs}}{\text{g\;sec}}\;
\left(\frac{C_A f_i f_e^2}{A}\right)
\left(\frac{\rho_0}{\text{g\;cm}^{-3}}\right)
T_K^{\,1/2}\bigl(1 + 4.4 \times 10^{-10}\,T_K\bigr)\,\bar{G}(T).
\tag{2.1.31}
\end{align}

\noindent (\citealt{Ginzburg1967}).


%%---------------------------------------------------------------------------
%%  SECTION 2.2 — Free-Bound Radiation
%%---------------------------------------------------------------------------

\subsection{Free-Bound Radiation}
\label{sec:2.2}

\textit{Basic references}\quad \citet{Jackson1962}; \citet{Ginzburg1967}.

\bigskip
\textit{Basic physics and formulas}

When an electron of kinetic energy
\begin{equation}
\tfrac{1}{2}m_e v^2 \leq Z^2 R_y
\end{equation}
impinges on an ion, there is a significant probability that, instead of scattering
with the emission of bremsstrahlung, it will get captured into a bound state with
an accompanying emission of ``free-bound'' radiation.

We shall not compute here the details of the radiation.  Instead we cite the
results of such a computation from the above references.

Consider a plasma at temperature $T$ and density of rest mass $\rho_0$.  Let $C_A$ be
the concentration by mass of a particular atom in the plasma; let $f_i$ be the fraction
of all such atoms ionized into a given state, with charge $Ze$; and let $n_q$ be the
principal quantum number of that state, and $E_i$ its
ionization energy.  Then the emissivity due to such transitions is
\begin{align}
\varepsilon_\nu &= 8c^2
\left(\frac{C_A f_i f_e^2}{An_q^5}\right)
\left(\frac{\alpha^2 m_e e^2}{m_p}\right)^{\!3/2}
\rho_0
\left(\frac{\rho_0}{3kT}\right)^{\!3/2}
e^{-(h\nu - E_i)/kT}\;G_{bf}
\notag\\[4pt]
&= 3.9 \times 10^{15}\;
\frac{\text{ergs}}{\text{g\;sec\;Hz}}\;
\left(\frac{C_A f_i f_e^2}{An_q^5}\right)
\left(\frac{\rho_0}{\text{g\;cm}^{-3}}\right)
T_K^{-3/2}\;e^{-(h\nu - E_i)/kT}\;G_{bf}.
\tag{2.2.1}
\end{align}

Here $G_{bf}$ is a ``bound-free Gaunt factor'', analogous to the ``free-free'' Gaunt
factors of \S\ref{sec:2.1}.  It depends on the kinetic energy
\begin{equation}
\tfrac{1}{2}m_e v^2 = h\nu - E_i
\tag{2.2.2}
\end{equation}
of the electron that is captured, and on the structure of the bound state (which
one characterizes by its quantum numbers $n_q, l_q, \ldots$):
\begin{equation}
G_{bf} = G_{bf}(h\nu - E_i;\;n_q, l_q, \ldots).
\tag{2.2.3}
\end{equation}

Gaunt factors for various bound-free transitions of astrophysical interest are
tabulated and graphed by \citet{Karzas1961}.  Energy conservation
requires that all photons emitted have energies greater than $E_i$; consequently,
\begin{equation}
G_{bf} = 0 \quad \text{for} \quad h\nu - E_i < 0.
\tag{2.2.4}
\end{equation}

In general, $G_{bf}$ ``turns on'' discontinuously at $h\nu = E_i$, with an initial value
between ${\sim}\,0.5$ and ${\sim}\,1.5$:
\begin{equation}
G_{bf} \simeq 1 \quad \text{for} \quad 0 < h\nu - E_i \lesssim E_i.
\tag{2.2.5}
\end{equation}

For $h\nu - E_i \gtrsim Z^2 R_y$, $G_{bf}$ typically remains within an order of magnitude of
unity; as $h\nu - E_i \geq 10\,Z^2 R_y$, $G_{bf}$ falls rapidly.
See \citet{Karzas1961}.

The total emissivity, $e = \int \varepsilon_\nu\,d\nu$, due to a given free-bound transition is
\begin{align}
e &= 8\sqrt{3}\,c^2
\left(\frac{C_A f_i f_e^2}{An_q^5}\right)
\left(\frac{\alpha^2 m_e e^2}{m_p^2}\right)^{\!3/2}
\rho_0
\left(\frac{2\pi m_e c^2}{3kT}\right)^{\!1/2}
T_K^{-1/2}\;\bar{G}_{bf}
\notag\\[4pt]
&= 4.2 \times 10^{26}\;
\frac{\text{ergs}}{\text{g\;sec}}\;
\left(\frac{C_A f_i f_e^2}{An_q^5}\right)
\left(\frac{\rho_0}{\text{g\;cm}^{-3}}\right)
T_K^{-1/2}\;\bar{G}_{bf}
\tag{2.2.6}
\end{align}

where $\bar{G}_{bf}$, the mean Gaunt factor, depends only on temperature and on the
bound state, and is approximately equal to one.
\begin{equation}
\bar{G}_{bf} = \int_0^\infty G_{bf}(h\nu - E_i = xkT;\;n_q, l_q, \ldots)\;e^{-x}\,dx \simeq 1.
\tag{2.2.7}
\end{equation}

The total emissivity (2.2.6) due to a given free-bound transition, compared
to the total emissivity (2.1.29) of free-free radiation from the same type of ion,
is
\begin{equation}
\frac{\varepsilon_{fb}}{\varepsilon_{ff}}
\simeq \frac{Z^2}{n_q^5}\,
\left(\frac{8 \times 10^5}{T}\right).
\tag{2.2.8}
\end{equation}

\citet{Ginzburg1967} states that for astrophysical plasmas the ratio of total free-bound
radiation to total free-free radiation (summed over all states and all ions)
generally is less than
\begin{equation}
\frac{\varepsilon_{fb\,\text{total}}}{\varepsilon_{ff\,\text{total}}}
< \left(\frac{8 \times 10^5\,\text{K}}{T}\right).
\tag{2.2.9}
\end{equation}
